\subsection{The exponent from an elliptic curve}
\label{subsec:curve-exponents}
\hypertarget{source-curve-exponents}{}

The logarithmic differences and torus action used in the jet estimates arise from a single elliptic formal group. We construct them in coordinates in which the Frobenius congruences can be compared. Work first over a finite
unramified extension of \(\mathbb Q_p\), or the completion of its maximal
unramified extension, with Frobenius lift \(\sigma\), \(p\ge5\).
Write \(f\sim h\) for zero-constant series if \(\exp(f-h)\)
is an integral unit with constant one. The relation is additive,
is preserved by integral substitutions fixing zero, and permits
multiplication by a scalar in \(\mathbb Z_p\). For the latter
property expand a principal unit to that power by the integral
binomial coefficients. These statements continue to hold after
valued extension.

Let 
\[
  L(X)=\sum_{n\ge0}p^{-n}X^{p^{2n}},
\]
 and let \(F\) be
the formal law with logarithm \(L\). We need precise control
of which additive-coordinate substitutions are integral in
formal group coordinates. On series in any finite set of
variables let \(\Phi\) apply \(\sigma^2\) to the coefficients
and raise the variables to \(p^2\). For \(A(0)=0\),
\begin{equation}\label{eq:integral-log}
 L^{-1}(A)\text{ is integral}\quad\Longleftrightarrow\quad
                       A-p^{-1}\Phi A\text{ is integral}.
\end{equation}
\begin{proposition}[The integral logarithm criterion]\label{prop:quadratic-integral-log}
On the unramified ring and for \(A(0)=0\), the condition
\eqref{eq:integral-log} characterizes integral substitutions into
\(L^{-1}\). In particular the following formal sum and scalar substitutions have coefficients in the unramified ring, with \(a\) integral in the scalar case:
\[
 F(X,Y)=L^{-1}(L(X)+L(Y)),\qquad [a]X=L^{-1}(aL(X))\quad(\sigma^2a=a).
\]

\end{proposition}

\begin{proof}
For \(A=L(h)\) the latter difference is \(h\) plus
\[
 \sum_{a\ge0}p^{-a-1}\big\{(h^{p^2})^{p^{2a}}-
                                      (\Phi h)^{p^{2a}}\big\}.
\]
If \(h\) is integral, then \(h^{p^2}\equiv\Phi h\bmod p\).
Raising a congruence modulo \(p^b,\ b\ge1\), to the \(p\)-th
power improves it to one modulo \(p^{b+1}\). This proves
integrality termwise. Conversely determine \(h\) by degree:
at each degree the terms in braces depend only on earlier
coefficients, so the same argument applied to that integral
truncation proves the next coefficient integral. In particular
\(F\) and \([a]X=L^{-1}(aL(X))\) for \(a\) integral,
\(\sigma^2a=a\), are integral.

\end{proof}

Use a short Weierstrass curve \(y^2=4x^3-g_2x-g_3\) with
integral invariants and good reduction. Let \(z\) be its
additive uniformizing parameter, \(x=\wp(z),y=\wp'(z)\),
and \(t=-2x/y\) the formal parameter. Then \(z=l(t)=t+\cdots\)
is its formal logarithm; \(x(t)=t^{-2}a(t)\) with \(a\)
integral, even, \(a(0)=1\). Indeed, substituting \(y=-2x/t\) and \(x=t^{-2}a\) in the curve
equation gives
\[
 P(t,a)=a^3-a^2-(g_2/4)at^4-(g_3/4)t^6=0,\qquad
 P(0,1)=0,\quad \partial_aP(0,1)=3-2=1 .
\]
The formal implicit-function theorem therefore gives a unique integral
solution with \(a(0)=1\).

Since \(P(-t,a)=P(t,a)\), uniqueness also
gives \(a(-t)=a(t)\). Write \(\varsigma\) for
the odd Weierstrass sigma function normalized by
\(\varsigma'(0)=1\). Its formal series is specified by
\(-(\log\varsigma)''=\wp\) and oddness with the usual
zero quadratic term in \(\log(\varsigma(z)/z)\).
For a scalar \(s_2\) put
\[
       q_E(t)=-\log\frac{\varsigma(l(t))}{t}
                                     +\tfrac12s_2 l(t)^2.
\]
Let \(t_\pm\) represent addition and subtraction on the curve.
\begin{proposition}[Additivity and conjugate scalars in the curve pairing]\label{prop:curve-additive-pairing}
Suppose the good-model change \(l^{-1}(cL(X))\) is integral
with \(c\) a unit. The parallelogram relation for \(q\) holds
modulo \(\sim\), and its polarization \(b\) is additive in both
variables modulo \(\sim\). If \eqref{eq:cm-norm-relation} holds
for all quadratic units, then it holds for all quadratic integers
and for each such \(a\) the polarized scalar action satisfies
\[
 b([a]x,y)\sim b(x,[\sigma(a)]y).
\]

\end{proposition}

\begin{proof}
\proofstep{The integral parallelogram relation.}

The formal addition formula for sigma gives
\[
 \frac{\varsigma(l(t)+l(u))\varsigma(l(t)-l(u))}
      {\varsigma(l(t))^2\varsigma(l(u))^2}=\wp(l(u))-\wp(l(t)).
\]
For instance, as elliptic functions of \(l(t)\) the two sides
have the same poles and principal parts by quasi-periodicity
and the normalization at zero. Their difference is constant,
and evaluation at \(l(t)=l(u)\) fixes that constant. Expansion
at the identity yields the formal identity over the field of
the invariants by the coefficient recursions for \(\wp\) and
\(\varsigma(z)/z\).

It follows that the exponential of minus the parallelogram
difference \(q_E(t_+)+q_E(t_-)-2q_E(t)-2q_E(u)\) is
\[
 \frac{t^2a(u)-u^2a(t)}{t^2-u^2}
                   \frac{t+u}{t_+}\frac{t-u}{t_-}.
\]
Each factor is an integral unit with constant one. Evenness
of \(a\) gives this for the first; integrality of the curve's
formal law and its odd inverse give divisibility by \(t\pm u\)
and the linear term gives the other two. For a unit \(c\)
suppose that \(t(X)=l^{-1}(cL(X))\) is integral. This is the
coordinate comparison required of the local CM model. Its
linear term is \(cX\). The exponent in these coordinates is
\[
           q(X)=-\log\frac{\varsigma(cL(X))}{cX}
                                    +\tfrac12s_2(cL(X))^2 .
\]
It differs from \(q_E(t(X))\) by a logarithm of an integral unit
since \(t(X)/(cX)\) is such a unit. The same parallelogram
relation therefore holds for \(q\) modulo \(\sim\).

\proofstep{Additivity via the integral formal half.}

Let \(b(X,Y)=q(X+_FY)-q(X)-q(Y)\). It vanishes on the two
axes and is symmetric. Modulo \(\sim\) it is additive in each
variable on all integral formal points. To see this, subtract
the parallelogram identities at \((x+z,y)\) and \((x,y)\),
writing the group law additively. For \(f(x)=b(x,z)\) this gives
\(f(x+y)+f(x-y)\sim2f(x)\). Taking \(y=x\) gives

\[
  f(2x)\sim2f(x)-f(0)=2f(x),
\]
 since \(f(0)=0\). With \(x=(u+v)/2\) and \(y=(u-v)/2\) in the formal group, the two
relations give
\[
 f(u)+f(v)\sim 2f\bigl((u+v)/2\bigr)\sim f(u+v).
\]
The formal half is integral here. Symmetry of \(b\) gives additivity
in the second variable.

\proofstep{Conjugate scalars and extension from units.}

There is also a scalar comparison from CM. Its needed input
on this model is that for every unit \(a\) fixed by \(\sigma^2\),
\begin{equation}\label{eq:cm-norm-relation}
                   q([a]X)\sim a\sigma(a)\,q(X).
\end{equation}
The scalar \(n=a\sigma(a)\) is in \(\mathbb Z_p^\times\).
Polarization gives \(b([a]x,[a]y)\sim n b(x,y)\).
For \(n\in\mathbb Z_p\), additivity implies
\(b(x,[n]y)\sim n b(x,y)\): approximate \(n\) by nonnegative
integers. Each coefficient of each side and its exponential
difference is continuous in \(n\), since at bounded degree
only finitely many operations occur. The valuation ring is
closed. Replacing \(y\) by \([a^{-1}]y\) therefore gives
\[
                        b([a]x,y)\sim b(x,[\sigma(a)]y).
\]
This holds for nonunits \(a\) fixed by \(\sigma^2\) as well,
by subtracting the identities for \(1+a\) and \(1\).

In particular \(q([a]x)\sim a\sigma(a)q(x)\) for nonunits too:
use \([a]x=[1+a]x-x\), the evenness of \(q\),
and \(b(x,x)\sim2q(x)\) from its parallelogram identity.
Moving \(1+a\) across this pairing, or polarizing
\(q([1+a]x)\), gives the claimed identity. Explicitly
\(b([a]x,x)+b(x,[a]x)\sim
b(x,[a+\sigma(a)]x)\sim2(a+\sigma(a))q(x)\);
this determines the cross term by division by two.

\end{proof}

Write \(\widetilde b(G,H)=b(L^{-1}G,L^{-1}H)\).
The preceding rules apply whenever \(G,H\) satisfy
\eqref{eq:integral-log}. For an even cycle choose integral
bases \(\omega_r,\omega_r^*\) of the unramified degree \(d\)
field over its unramified quadratic subfield, dual for the
relative trace. Such bases exist since the residue trace is
perfect and its lifted determinant is a unit. Set
\[
 \begin{aligned}
 E(c,Z)&=\sum_{s\ {\rm even}}\sigma^s(c)Z_s,&
 O(c,Z)&=\sum_{s\ {\rm odd}}\sigma^s(c)Z_s,\\
 B^\circ(Z)&=\sum_r\widetilde b(E(\omega_r^*,Z),O(\omega_r,Z)).
 \end{aligned}
\]
Define \(B\) as the average of \(B^\circ\) under
\(Z_s\mapsto u^{(-1)^s p^s}Z_s\) over \(u^{p^d-1}=1\). It
has coefficients in an unramified field and is torus invariant.
The two positive parity degrees follow from the axes of \(b\).
An even-odd bilinear term would have character
\(p^s-p^t\bmod(p^d-1)\) for distinct \(s,t\), so the averaging
eliminates it.

\begin{proposition}[The cycle exponent in a trace-dual frame]\label{prop:curve-cycle-exponent}
On an even cycle, with the integral curve and scalar comparisons
just proved, the averaged exponent \(B\) has the two integrality
properties of Lemma~\ref{lem:two-differences}. The trace-dual
substitution from an integral orthogonal theta chart gives
the local series of Proposition~\ref{prop:cyclic-jet}, whose exponential factor and integral remainder combine to give
\[
 \exp\!\left(-\sum_{i,b}B((Z_{sib})_s)\right)U(g(Z)).
\]
 The one-letter exponent is the equal-variable
specialization with the factor \(1/2\).
\end{proposition}

\begin{proof}
\proofstep{Check the two parity differences.}

For \(Z_s={\cal L}_s(X)\) one has
\(Z_s-p^{-1}\Phi Z_{s-2}=X_s\). More generally suppose all
these differences are integral and all \(Z_s(0)=0\). Both
\(E(c,Z)\) and \(O(c,Z)\) satisfy the right side of
\eqref{eq:integral-log} for integral \(c\) fixed by \(\sigma^d\),
by shifting indices in the sum. Torus translation replaces
\(\omega_r^*\) by \(u\omega_r^*\) and \(\omega_r\) by
\(u^{-1}\omega_r\), since 
\[
  \sigma^s(u)=u^{p^s}.
\]
The parity differences of each translated \(B^\circ\), evaluated
on the two copies of \({\cal L}\), are therefore equivalent to
zero by additivity. This remains true on averaging since
\((p^d-1)^{-1}\in\mathbb Z_p\). This verifies precisely the two
integrality hypotheses of Lemma~\ref{lem:two-differences}.

\proofstep{Compare the torus average to the original contraction.}

To identify the resulting exponent with the one on the original
theta chart, one needs \(B(Z)\sim B^\circ(Z)\) on these points.
For a torus element write
\(c_{rr'}=\operatorname{Tr}(u\omega_r^*\omega_{r'})\), with
the trace over the quadratic subfield. Its entries are integral,
fixed by \(\sigma^2\). Then
\[
 u\omega_r^*=\sum_{r'}c_{rr'}\omega_{r'}^*,\qquad
 \omega_{r'}=\sum_r c_{rr'}\,u^{-1}\omega_r .
\]
In \(E\) the coefficients of the first expansion are \(c_{rr'}\).
Move each across \(\widetilde b\) by the scalar comparison;
the resulting \(\sigma(c_{rr'})\) are exactly the coefficients
in \(O\) for the second expansion. This proves invariance
modulo \(\sim\) under each torus element, whence the asserted
comparison with the mean.

\proofstep{Apply the hyperbolic Gram identities.}

For clarity, suppose the theta series on the orthogonal chart
has \(U_0=\exp(\sum_i q(X_i))\Theta(X)\) integral. To pass to
a hyperbolic frame use the formal sum of multiples of independent
coordinates with coefficients \(f_{ei},f'_{ei}\) integral in
the quadratic subfield. The orthogonality conditions on these two families are
\[
  \sum_i\sigma(f_{ei})f_{e'i}=0,\qquad\sum_i\sigma(f'_{ei})f'_{e'i}=0,\qquad\sum_i\sigma(f_{ei})f'_{e'i}=\delta_{ee'}.
\]
Expanding \(q\) on this formal sum gives diagonal terms
\(f_{ei}\sigma(f_{ei})q(X_e)\) and terms paired by \(b\),
with the conjugated first coefficient moved to the second slot.
Summing over \(i\) leaves just \(\sum_e b(X_e,Y_e)\) modulo
\(\sim\). Thus after this substitution the theta series is an
integral series times the exponential of minus that sum.
The index \(e\) includes a trace-basis index and a matrix
entry.

Now set
\[
  X_r=L^{-1}(E(\omega_r^*,Z)),\qquad
  Y_r=L^{-1}(O(\omega_r,Z))
\]
at each entry. On substituting \(Z={\cal L}(W)\), both are
integral formal points by \eqref{eq:integral-log}. The
exponential of the difference between the summed exponents
and the sum of the corresponding \(B\)'s is an integral unit
in \(W\). Absorb it into the integral factor and substitute
\(W=g(Z)\). This gives the factor \(U(g(Z))\) and the exponent
\(-\sum_{i,b} B((Z_{sib})_s)\) of
Proposition~\ref{prop:cyclic-jet}. Auxiliary orbit variables
can be left as indeterminates throughout. For one letter use
\(d=2,\ \omega=\omega^*=1\). In the equal-variable
specialization 
\[
  b(X,X)=q([2]X)-2q(X)\sim 2q(X);
\]
 division by two yields
the one-letter exponent. A unit change of the logarithmic
variable common to the two slots scales the homogeneous jet by
a unit.

This calculation gives the local series from an integral
orthogonal theta chart, an integral logarithm comparison and
\eqref{eq:cm-norm-relation}. Establishing those three inputs
on the family requires its models and the normalization of its
translated theta sections.
\label{cor:curve-exponents-1}
\end{proof}
